\subsection{An explicit reduced counterexample to unrestricted C.4}

\begin{theorem}[C-RD10: infinitely many nodes obstruct the local-system lemma]
\label{thm:C-RD10}
\status{Proved; independently audited.}
There is a countable, one-dimensional, noetherian integral
$\mathbf Q$-algebra $R$, all of whose closed local rings are ordinary
split nodes, and a generization-compatible family of elliptic schemes
on $\operatorname{Spec}\mathcal O_{\operatorname{Spec}R,s}$ which does
not come from an abelian scheme over $\operatorname{Spec}R$.
Consequently C.4 is false for general reduced bases and hence for
arbitrary locally noetherian bases. The counterexample is explicit
and uses no cardinality diagonalization.
\end{theorem}
\begin{proof}
\emph{The ring.} Let $t_1,t_2,\ldots$ be algebraically independent over
$\mathbf Q$, set $x_i=t_i^2-1$ and $y_i=t_i(t_i^2-1)$, and put
\[
 C_i=\mathbf Q[x_i,y_i]
   \simeq\mathbf Q[X,Y]/(Y^2-X^2(X+1)),\qquad
 R_0=\bigotimes_{i\geq1,\mathbf Q} C_i\subset\mathbf Q[t_1,t_2,\ldots].
\]
Write $P_i=(x_i,y_i)R_0$, set
\[
 \Sigma=R_0\setminus\bigcup_{i\geq1}P_i,
 \qquad R=\Sigma^{-1}R_0,
 \qquad \mathfrak m_i=P_iR.
\]
Each $P_i$ is prime, since its quotient is the tensor product of the
other geometrically integral nodal rings. Thus $\Sigma$ is
multiplicative. A nonzero element of $R_0$ involves only finitely many
indices, and belongs to none of the $P_i$ outside those indices.

We verify all the finiteness properties used below. If a nonzero prime
$Q\subset R_0$ avoids $\Sigma$, choose $0\ne f\in Q$. The set
$I=\{i:f\in P_i\}$ is finite. If $Q$ were contained in none of these
$P_i$, finite prime avoidance would give
$g\in Q\setminus\bigcup_{i\in I}P_i$. Let $J$ contain every index
occurring in $f$ and $g$. For $i\in I$, $f+cg\notin P_i$ whenever
$0\ne c\in\mathbf Q$. For each $i\in J\setminus I$ there is at most
one scalar $c$ with $f+cg\in P_i$. Choose a nonzero rational number
avoiding these finitely many exceptions and also ensuring $f+cg\ne0$.
For $i\notin J$, independence of $t_i$ gives $f+cg\notin P_i$.
Then $f+cg\in Q\cap\Sigma$, a contradiction. Every prime of $R$
is therefore contained in some $\mathfrak m_i$. These ideals are
pairwise incomparable: $x_i\in P_i\setminus P_j$ for $i\ne j$, a
property unchanged by localization. Hence they are exactly the
maximal ideals of $R$.

Let $L_i=\mathbf Q(t_j:j\ne i)$. Localizing first in the nonzero
coefficients from the other factors gives
\[
 R_{\mathfrak m_i}
   \simeq L_i[X,Y]_{(X,Y)}/(Y^2-X^2(X+1)).                 \tag{C-RD10.1}
\]
These are noetherian one-dimensional local domains. Every nonzero
$a\in R$ belongs to only finitely many maximal ideals. If $I\subset R$
is a nonzero ideal, choose $0\ne a\in I$. At each of those finitely
many maximal ideals containing $a$, choose finitely many elements of
$I$ generating its localization. Together with $a$ they generate an
ideal $J\subset I$ whose localization equals $I$ at every maximal
ideal; therefore $I/J=0$. Thus every ideal is finitely generated,
proving that $R$ is noetherian. Formula (C-RD10.1) proves $\dim R=1$.
Its points are the generic point $\eta$ and the $\mathfrak m_i$.
The tangent cone $Y^2-X^2$ has two distinct rational branches.

The element $t_i=y_i/x_i$ lies in $F=\operatorname{Frac}R$ and satisfies
$t_i^2=1+x_i$, but does not belong to $R_{\mathfrak m_i}$. Indeed,
every rational function in (C-RD10.1) regular at the node has equal
values at the two points $t_i=1$ and $t_i=-1$ of its normalization:
write it as $a(X,Y)/b(X,Y)$ with $b(0,0)\ne0$ and evaluate. The
function $t_i$ has the distinct values $1$ and $-1$.
This also proves that the normalization of $R$ is not finite. If it
were finite, some nonzero $d\in R$ would multiply it into $R$; choose
$i$ with $d\notin\mathfrak m_i$. The integral element $t_i$ would
then belong to $R_{\mathfrak m_i}$, a contradiction.

\emph{The local elliptic schemes.} Fix
\[
 E/\mathbf Q:\quad Y^2=X^3+X+1,
 \qquad \Delta(E)=-496.
\]
At the generic point use $E_F$. At $\mathfrak m_i$, put $d_i=1+x_i$
and take
\[
 A(\mathfrak m_i):\quad
 Y^2=X^3+d_i^2X+d_i^3
 \quad\text{over }R_{\mathfrak m_i}.                    \tag{C-RD10.2}
\]
The element $d_i$ is a unit, and the discriminant is $-496d_i^6$,
so this is an elliptic scheme. Its generic fibre is identified with
$E_F$ by $(X,Y)\mapsto(X/t_i^2,Y/t_i^3)$. Since the only nonidentity
generizations on $\operatorname{Spec}R$ run from a closed point to
$\eta$, these identifications define the required local system.

However (C-RD10.2) is not isomorphic to the constant $E$ over
$R_{\mathfrak m_i}$. Since $2$ and $3$ are invertible, an isomorphism
of short Weierstrass elliptic schemes preserving the zero section is
of the form $X\mapsto u^2X$, $Y\mapsto u^3Y$, with $u$ a unit;
comparison of the absent $X^2$, $XY$ and $Y$ terms eliminates all
translations. Comparing the two nonzero coefficients gives
$u^4=d_i^2$ and $u^6=d_i^3$, hence $u^2=d_i$. In the field $F$ this
forces $u=\pm t_i$, which are not elements of $R_{\mathfrak m_i}$.

\emph{Failure of effectivity.} If an abelian scheme $A/R$ realized the
system, its generic fibre would be identified with $E_F$. Finite
presentation spreads this generic isomorphism and its inverse to
some principal dense open $D(f)$, with $0\ne f\in R$. But $f$ lies
in only finitely many $\mathfrak m_i$. For every other $i$, one would
have
\[
 A(\mathfrak m_i)\simeq A\times_RR_{\mathfrak m_i}
   \simeq E\times_{\mathbf Q}R_{\mathfrak m_i},
\]
contradicting the preceding paragraph. This proves noneffectivity. The exact polynomial identities and
Weierstrass-change coefficients are also verified in
\texttt{verify\_nodal\_twist.py}, with its saved JSON certificate.
\end{proof}

\begin{remark}[Scope of the example]
The infinite tensor-product localization follows the construction
attributed to \cite[Proposition~1 and Example~1]{Hochster1973Loci},
with a nodal factor in place of the cuspidal one. The original
Hochster article has not yet been obtained: author bibliographical
metadata and the Stacks reference were checked, and the proposition
was located in a secondary quotation. The preceding proof independently
checks noetherianity and every property used in the counterexample.
It does not contradict Theorem~\ref{thm:C-RD8}: its conclusion in fact
shows that $\operatorname{Spec}R$ has no open cover by schemes admitting
finite surjective geometrically unibranch covers. The single-node
$[-1]$ proposal in Lead C remains effective; the obstruction here is
that infinitely many nontrivial local twists cannot be supported on
the complement of any dense open.
\end{remark}
